% =============================================================================
% EINLEITUNG ZU BAND 8: STANDARDMODELL AUS GF(27), LEPTONMASSEN UND
% SCHNITTSTELLEN (Bd. 8 von 8)
% =============================================================================

\section*{Achter Band der Dokumentensammlung}

Band~8 umfasst die Dokumente 344 bis 389 aus der Zeit von September bis
Anfang Oktober 2026 und schließt die Sammlung auf dem aktuellen Stand des
Korpus ab. Er baut unmittelbar auf der Galois-Brücke aus Band~7 auf und
führt sie bis zu den Quantenzahlen des Standardmodells, zu den Leptonmassen
und zu den Schnittstellen mit anderen theoretischen Rahmen weiter.

\subsection*{Inhalt und Schwerpunkt}

Der erste Teil leitet aus der Struktur von $GF(27)^*$ die
Ladungsquantisierung, die Gell-Mann--Nishijima-Relation
$Q=I_3+Y/2$, die Generationenstruktur mit ihrer Kopplungsmatrix und die
$SU(2)_L$-Chiralität ab; ergänzt wird dies durch den Vergleich mit den
Furey-Fermionen, eine Einordnung in die Literatur und die Behandlung
Schwarzer Löcher unter der Zeit-Masse-Dualität. Der zweite Teil klärt
zunächst, wie die PDG-Vergleichswerte zustande kommen und welche
Modellannahmen in sie eingehen, und stellt dann Leptonleiter, Galois-Ansatz
und Koide-Relation gegenüber; es folgen die Koide-Amplitude $d/c=\sqrt2$, die
Higgs-Geometrie, die spontane Symmetriebrechung aus der
$T^4/\mathbb{Z}_3$-Geometrie und die Übersicht der Standardmodell-Teilchen
mit ihrem Belegstatus.

Der dritte Teil zeigt, warum Harmonik und algebraische Struktur dieselbe
Mathematik sind, und prüft diese Einsicht an Laborsystemen: Frequenzanalyse
und ihre Grenzen, HRV-Messung, Kagome-Gitter und die BAW-Ising-Maschine.
Hinzu kommen die Hodge-Theorie auf $T^4/\mathbb{Z}_3$, der Vergleich von
Ikosaeder-Träger und $D_4$-Gitter, die Darstellung des Fundaments in vier
Schichten und die Feldstruktur von Energie, Fluss und Geometrie. Der vierte
Teil ist den Leptonmassen gewidmet: $\theta=p_0=2/9$ als zwei Lesarten eines
Quotienten, der Weg von $p_0,p_1,p_2$ zu den Massen, vier Wege im Vergleich,
die Begründung des Ikosaeders, der Faktor $4/3$, das Matrixelement-Prinzip
und der Yukawa-Mechanismus als Folge von $\tilde T\cdot m=1$.

Der fünfte Teil sammelt die Schnittstellen: die Brückengleichungen zu
Observer Patch Holography, die Hierarchie $v/E_P$, eine Übersicht aller
Vergleichs- und Brückendokumente, Überlegungen zur Begutachtung, die
Spezifität des $D_4$-Trägers, die exakt aufsummierte Rekursion, den Stand
des Myon-$g-2$, die Massen ohne $v$, eine Kurzfassung der FFGFT nach der
Rechenprüfung, den Higgs-Vakuum-Weg zu $\xi$, die Lage von $\alpha$ in der
Geometrie, eine Übersicht der Abweichungen, die CMB-Temperatur in der
Grundform und die Darstellungen der Gravitation am Beispiel des
Parabelflugs.

\subsection*{Gliederung}

Die 46 Dokumente sind in fünf Teile gegliedert: Quantenzahlen aus der
Galois-Struktur (Dok.~344--350); Vergleichswerte, Leptonreste, Higgs und
Standardmodell-Teilchen (Dok.~351--357); Harmonik, Laborsysteme,
Hodge-Theorie und Felder (Dok.~358--366); Ikosaeder, Matrixelemente und
Leptonmassen (Dok.~367--373); Schnittstellen, Begutachtung und Grundform
(Dok.~374--389).

\subsection*{Zum Lesen der Dokumente}

Für einen raschen Überblick über den Gesamtstand eignet sich Dok.~384
(Kurzfassung), für die verbleibenden Abweichungen Dok.~387. Statusmarker
und Korrekturregister gelten wie in den vorangehenden Bänden; maßgeblich ist
jeweils die jüngste Fassung eines Ergebnisses im Korpus.

\subsection*{Stand}

Die Dokumente dieses Bandes tragen den Stand ihrer jeweiligen Fassung im
Korpus zum Oktober 2026.
